\begin{figure} [h]
	\centering
	\resizebox{\textwidth}{!}{	
		\tikzsetnextfilename{fig_short_perturbations_implemented}		
		\begin{tikzpicture}
			\centering
			%
			\begin{axis}[%
				name=ax1,
				width=.48\linewidth,
				height=.2\linewidth,
				scale only axis,
				xlabel={temps (\si{\milli\second})},
				ylabel={Force (\si{\newton})},
				axis x line = bottom,
				axis y line = left,
				ylabel near ticks,
				]
				%
				\addplot [
				color=blue,
				smooth,
				]
				table [col sep=comma, x expr=\coordindex, y=force_pert_pos] {misc/short_force_perturbations.csv};
			\end{axis}
			%
			\begin{axis}[%
				width=.48\linewidth,
				height=.2\linewidth,
				scale only axis,
				xlabel={temps (\SI{}{\milli\second})},
				%ylabel={Force (\SI{}{\newton})},
				axis x line = bottom,
				axis y line = left, 			
				yticklabel style={
					%/pgf/number format/precision=3,
					/pgf/number format/fixed},
				at={($(ax1.south east)+(6,0)$)},
				%scaled y ticks=false,
				]
				%
				\addplot+[
				%color=blue,
				smooth,
				mark=none,
				]
				table [col sep=comma, x expr=\coordindex, y=force_pert_neg] {misc/short_force_perturbations.csv};
			\end{axis}
			
		\end{tikzpicture}
	}
	\caption{Exemples d'une perturbation positive et d'une négative, qui ont pu être extraites en utilisant la méthode décrite en \cref{chap3_sec_method_sub_virt_traj_6}.}
	\label{fig_short_perturbations_implemented} % Label 
\end{figure}